% !PRINTED none
\weberpage{none}

\begin{center}
\gesperrt{Tenth Chapter.}

\bigskip

\gesperrt{General Theory of Metacyclic Equations.}
\end{center}

\bigskip
\begin{center}
\rule{4cm}{0.4pt}
\end{center}
\bigskip

\begin{center}
\S.\,73.

\medskip

\gesperrt{The Resolvents of the Composition Series.}
\end{center}

\bigskip

From the theorems of general group theory that are treated in the first Book of this
volume, important algebraic consequences result when they are applied to the
\gesperrt{Galois} group of an equation.

An arbitrary field $\Omega$ is now assumed as domain of rationality; let
$f(x) = 0$ be an equation in this field without multiple roots, and $P$ its
\gesperrt{Galois} group. Let this group $P$ have a normal divisor $Q$. We denote by
$n$ the degree of $P$ and by $j$ the index of the divisor $Q$, so that $n : j$ is the
degree of $Q$.

We have seen in the first volume (\S.\,156) that the group of $f(x)$ is reduced to
$Q$ by the adjunction of a root of an irreducible normal equation of degree $j$,
and this auxiliary equation of degree $j$ we have designated
the \gesperrt{partial resolvent}.

The \gesperrt{Galois} group of this partial resolvent we obtained as follows: if
$\psi$ denotes a function of the roots of $f(x)$ belonging to the group $Q$, and if $P$
is decomposed into the subsidiary groups
\[
P = Q + Qa + Qb + Qc + \cdots
\]
where, then, $a, b, c, \ldots$ are certain permutations from $P$, then $\psi$ passes,
under the whole subsidiary group $Qa$, into one and the same function $\psi_a$, and the
\gesperrt{Galois} group of the resolvent consists of the substitutions
\[
(\psi, \psi), (\psi, \psi_a), (\psi, \psi_b), (\psi, \psi_c), \ldots
\]
